% =====================================================================
%  PREAMBLE DÙNG CHUNG (dựa trên mẫu Lab1--Lab4 của nhóm)
%  File gốc phải khai báo TRƯỚC khi \input file này:
%     \newif\ifsolution  \solutiontrue / \solutionfalse
% =====================================================================
\usepackage{iftex}
\ifPDFTeX
  \usepackage[utf8]{inputenc}
  \usepackage[T5]{fontenc}
  \usepackage{lmodern}
\else
  \usepackage{fontspec}
\fi
\usepackage[vietnamese,english]{babel}   % ngôn ngữ chính: English; báo cáo gọi \selectlanguage{vietnamese}
\usepackage[top=2cm,bottom=2cm,left=2cm,right=2cm]{geometry}
\usepackage{amsmath,amssymb}
\usepackage{graphicx}
\usepackage{tikz}
\usepackage{float}
\usepackage{caption}
\usepackage{enumitem}
\usepackage{fancyhdr}
\usepackage{lastpage}
\usepackage{booktabs}
\usepackage{colortbl}
\usepackage{array}
\usepackage{tabularx}
\usepackage{longtable}
\usepackage{titlesec}
\usepackage{comment}
\usepackage[colorlinks=true,urlcolor=blue,linkcolor=black]{hyperref}
\usepackage[most]{tcolorbox}
\usepackage{needspace}
\usepackage{tocloft}
% Mục lục 3 cấp: Lab N  >  N.k phần lớn  >  mục nhỏ / từng bài tập
\renewcommand{\cftchappresnum}{Lab~}
\setlength{\cftchapnumwidth}{3.6em}
\setlength{\cftsecnumwidth}{2.6em}
\setlength{\cftsubsecindent}{3.8em}
\setlength{\cftsubsecnumwidth}{4.9em}
\renewcommand{\cftsubsecfont}{\small}
\renewcommand{\cftsubsecpagefont}{\small}

\captionsetup{font=it,labelfont=it}

% ---------------------------------------------------------------------
%  Header / footer
% ---------------------------------------------------------------------
\pagestyle{fancy}
\fancyhf{}
\fancyfoot[L]{\itshape \courseNameEN\ (\courseCode)}
\fancyfoot[R]{\itshape Page \thepage}
\renewcommand{\headrulewidth}{0pt}
\renewcommand{\footrulewidth}{0.4pt}
\fancypagestyle{plain}{\fancyhf{}\fancyfoot[L]{\itshape \courseNameEN\ (\courseCode)}\fancyfoot[R]{\itshape Page \thepage}}

% ---------------------------------------------------------------------
%  Mỗi chương = 1 bài lab  ("Lab 1: ...")
% ---------------------------------------------------------------------
\titleformat{\chapter}[display]{\bfseries\huge\color{blue!60!black}}{\Large Lab \thechapter}{0.3em}{}[\vspace{2pt}{\color{blue!60!black}\titlerule[1pt]}]
\titlespacing*{\chapter}{0pt}{-1.2cm}{1.2em}
\setcounter{secnumdepth}{2}
\setcounter{tocdepth}{2}

% ---------------------------------------------------------------------
%  Đường dẫn ảnh của lab hiện tại:  \setlab{3}  ->  labs/lab3/images/...
% ---------------------------------------------------------------------
\newcommand{\labdir}{labs/lab0}
\newcommand{\setlab}[1]{\renewcommand{\labdir}{labs/lab#1}}

% Ô dán ảnh: tự hiện ảnh nếu tệp tồn tại, nếu chưa có thì hiện khung cam
\newcommand{\figslot}[4]{%
  \begin{figure}[H]
    \centering
    \IfFileExists{#1}{\includegraphics[#2]{#1}}{%
      \begin{tcolorbox}[enhanced,center,width=0.92\linewidth,height=#4,valign=center,halign=center,
          colback=yellow!8,colframe=orange!80!black,boxrule=1pt,arc=1.5mm]
        \textbf{IMAGE SLOT}\\[2pt]
        \edef\figpath{#1}%
        {\small\itshape Save the LTspice screenshot as\\ \texttt{\expandafter\detokenize\expandafter{\figpath}}\\ (this box disappears once the file exists)}
      \end{tcolorbox}}%
    \caption{#3}
  \end{figure}}

% \guidefig{tệp}{chú thích}      -> labs/labN/images/guide/tệp.png  (ảnh hướng dẫn thao tác)
% \circuitfig{tệp}{chú thích}    -> labs/labN/images/circuit/tệp.png (schematic chụp từ LTspice)
% \simfig[cao]{tệp}{chú thích}   -> labs/labN/images/sim/tệp.png    (đồ thị / kết quả)
% \installfig{tệp}{chú thích}    -> images/install/tệp.png
\newcommand{\guidefig}[2]{\figslot{\labdir/images/guide/#1.png}{width=0.85\linewidth,height=0.36\textheight,keepaspectratio}{#2}{4.5cm}}
\newcommand{\circuitfig}[2]{\figslot{\labdir/images/circuit/#1.png}{width=0.75\linewidth,height=0.32\textheight,keepaspectratio}{#2}{5cm}}
\newcommand{\simfig}[3][6cm]{\figslot{\labdir/images/sim/#2.png}{width=0.94\linewidth,height=0.42\textheight,keepaspectratio}{#3}{#1}}
\newcommand{\installfig}[2]{\figslot{images/install/#1.png}{width=0.85\linewidth,height=0.36\textheight,keepaspectratio}{#2}{4.5cm}}

% ---------------------------------------------------------------------
%  BÀI TẬP
% ---------------------------------------------------------------------
\newcounter{exercise}[chapter]
\newcounter{advexercise}[chapter]
\newcommand{\twodigit}[1]{\ifnum\value{#1}<10 0\fi\arabic{#1}}

% Nhãn trạng thái: Keep / Adjust / Replace
\newcommand{\statusbadge}[1]{%
  \def\st{#1}\def\stK{Keep}\def\stA{Adjust}\def\stR{Replace}%
  \ifx\st\stK\def\stcol{green!45!black}\else
  \ifx\st\stA\def\stcol{orange!85!black}\else
  \ifx\st\stR\def\stcol{red!70!black}\else\def\stcol{gray}\fi\fi\fi
  \tcbox[on line,colback=\stcol,colframe=\stcol,coltext=white,boxsep=1pt,left=4pt,right=4pt,top=1pt,bottom=1pt,arc=1mm]{\small\bfseries #1}}

% \exercise{Tiêu đề}{Keep|Adjust|Replace}{Bài gốc}
\newcommand{\exercise}[3]{%
  \ifnum\value{exercise}>0 \clearpage\else\bigskip\fi   % bài 1 nằm ngay dưới mục N.3
  \refstepcounter{exercise}%
  \subsection*{Exercise \thechapter.\arabic{exercise}: #1}%
  \addcontentsline{toc}{subsection}{\protect\numberline{Ex~\thechapter.\arabic{exercise}}#1}%
  \noindent\statusbadge{#2}\hfill{\small\itshape Original: #3 \quad|\quad LTspice file: \texttt{Lab\thechapter/Ex\twodigit{exercise}.asc}}\par\medskip}

% \advexercise{Tiêu đề}{Kỹ thuật chính (.step, .meas, .temp, Monte Carlo, .four ...)}
\newcommand{\advexercise}[2]{%
  \ifnum\value{advexercise}>0 \clearpage\else\bigskip\fi
  \refstepcounter{advexercise}%
  \subsection*{Advanced Exercise \thechapter.A\arabic{advexercise}: #1}%
  \addcontentsline{toc}{subsection}{\protect\numberline{Adv~\thechapter.A\arabic{advexercise}}#1}%
  \noindent\tcbox[on line,colback=violet!70!black,colframe=violet!70!black,coltext=white,boxsep=1pt,left=4pt,right=4pt,top=1pt,bottom=1pt,arc=1mm]{\small\bfseries Advanced}%
  \hfill{\small\itshape Technique: #2 \quad|\quad LTspice file: \texttt{Lab\thechapter/Adv\twodigit{advexercise}.asc}}\par\medskip}

% Lý do thay đổi (chỉ in trong bản lời giải)
\newcommand{\statusreason}[1]{\ifsolution\begin{tcolorbox}[colback=orange!6,colframe=orange!70!black,boxrule=0.6pt,arc=1mm,left=6pt,right=6pt,top=3pt,bottom=3pt,fonttitle=\bfseries\small,title=Reason for status (instructor only)]\small #1\end{tcolorbox}\fi}

% Đề bài
\newcommand{\task}[1]{\noindent #1\par\medskip}

% Linh kiện & mô hình
\newcommand{\components}[1]{\begin{tcolorbox}[colback=blue!3,colframe=blue!50!black,boxrule=0.6pt,arc=1mm,left=6pt,right=6pt,top=3pt,bottom=3pt,fonttitle=\bfseries\small,title=Components and models]\small #1\end{tcolorbox}}

% ---------------------------------------------------------------------
%  PHẦN DÀNH CHO SINH VIÊN / GIẢNG VIÊN
% ---------------------------------------------------------------------
\newcommand{\dotline}[1]{\noindent #1\dotfill\par\vspace{0.8em}}
\newcommand{\dotlines}[1]{\foreach \i in {1,...,#1}{\noindent\dotfill\par\vspace{1em}}}

% \answerspace{n}: n dòng chấm để sinh viên điền (ẩn trong bản lời giải)
\newcommand{\answerspace}[1]{\ifsolution\else\dotlines{#1}\fi}

% \ans{giá trị}: bản SV in ô trống, bản lời giải in giá trị
\newcommand{\fillcell}{\fcolorbox{orange!80!black}{yellow!25}{\rule{0pt}{1.1em}\rule{1.3cm}{0pt}}}
\newcommand{\ans}[1]{\ifsolution #1\else\fillcell\fi}

% Môi trường lời giải: chỉ xuất hiện khi \solutiontrue
\newtcolorbox{solbox}{enhanced,breakable,colback=blue!3!white,colframe=blue!60!black,fonttitle=\bfseries,
  title=Solution,arc=1.5mm,boxrule=0.8pt,left=8pt,right=8pt,top=6pt,bottom=6pt}
\ifsolution
  \newenvironment{solution}{\begin{solbox}}{\end{solbox}}
\else
  \excludecomment{solution}
\fi

% Bảng so sánh lý thuyết -- mô phỏng
%   \begin{comparison}{chú thích} \cmprow{Đại lượng}{Đơn vị}{LT}{MP}{Sai số} ... \end{comparison}
\newenvironment{comparison}[1]{%
  \def\cmpcaption{#1}%
  \begin{table}[H]\centering\renewcommand{\arraystretch}{1.3}%
  \begin{tabular}{l c c c c}\toprule
  \textbf{Quantity} & \textbf{Unit} & \textbf{Theory} & \textbf{Simulation} & \textbf{Error (\%)}\\\midrule}
  {\bottomrule\end{tabular}\caption{\cmpcaption}\end{table}}
\newcommand{\cmprow}[5]{#1 & #2 & \ans{#3} & \ans{#4} & \ans{#5}\\}

% Ô ghi nhận xét
\newtcolorbox{notebox}[2][]{enhanced,breakable,colback=gray!7,colframe=gray!65!black,fonttitle=\bfseries,
  title={#2},boxrule=0.8pt,arc=1.5mm,left=8pt,right=8pt,top=6pt,bottom=6pt,#1}

% Mẹo LTspice
\newtcolorbox{tipbox}{enhanced,breakable,colback=green!4,colframe=green!45!black,fonttitle=\bfseries,
  title=LTspice tip,boxrule=0.6pt,arc=1.5mm,left=8pt,right=8pt,top=4pt,bottom=4pt}

% Khối lệnh SPICE
\newtcblisting{spice}[1][]{listing only,colback=gray!5,colframe=gray!60!black,boxrule=0.6pt,arc=1mm,
  title=SPICE directives,fonttitle=\bfseries\small,
  listing options={basicstyle=\ttfamily\small,columns=fullflexible,keepspaces=true},#1}

% Ghi chú TODO cho nhóm (đổi \showtodotrue -> \showtodofalse khi nộp)
\newif\ifshowtodo \showtodotrue
\newcommand{\todo}[1]{\ifshowtodo{\color{red!70!black}\small\itshape [TODO: #1]}\fi}

% ---------------------------------------------------------------------
%  YÊU CẦU ĐẦU RA (Output Requirements) -- shared/output_common.tex
%   \begin{outputlist} \outrow{Ex~1}{Tiêu đề}{Đầu ra bắt buộc} ... \end{outputlist}
% ---------------------------------------------------------------------
\newenvironment{outputlist}{%
  \small\renewcommand{\arraystretch}{1.3}%
  \begin{longtable}{>{\centering\arraybackslash}p{1.5cm} >{\raggedright\arraybackslash}p{4.3cm} >{\raggedright\arraybackslash}p{8.4cm} >{\centering\arraybackslash}p{0.8cm}}
  \toprule
  \textbf{No.} & \textbf{Exercise} & \textbf{Required outputs in the report} & \textbf{Done}\\ \midrule\endhead
  \bottomrule\endfoot}
  {\end{longtable}}
\newcommand{\outrow}[3]{#1 & #2 & #3 & $\square$\\}
